\par\Needspace{9\baselineskip}\bigskip
\pdfbookmark[3]{Exposition, [23]}{pallas-II-article23}
\begin{center}[23.]\end{center}
\begingroup
\setlength{\lineskip}{4pt}\setlength{\lineskiplimit}{4pt}
\noindent\hspace*{\parindent}Par cette méthode on déterminera les élémens \(\AscendingNode+\delta\AscendingNode,\ i+\delta i,\ \varphi+\delta\varphi,\ \tilde\omega+\delta\tilde\omega,\ N+\delta N,\ n+\delta n\). Dans le développement de la variation instantanée de \(n+\delta n\) on s’aperçoit que les coëfficiens de tous les termes qui ont \(i=0\) deviennent égaux à zéro, d’où il s’ensuit que ni le mouvement moyen ni le grand axe ne contient une variation séculaire proprement dite, du moins en s’arrêtant à la première puissance de la force perturbatrice: aussi ces élémens ne contiendront pas de termes périodiques indépendants de \(L\). Pour le prouver d’une manière directe, reprenons la formule \(m\mathrm{d}\dfrac{1}{a}\bigl(\text{ou plutôt }m\mathrm{d}\dfrac{1}{a+\delta a}\bigr)=2(\xi\mathrm{d}x+\eta\mathrm{d}y+\zeta\mathrm{d}z)\). On voit facilement qu’en faisant
\[
m\left\{\frac{xx'+yy'+zz'}{r'^{8}}-\frac{1}{\rho}\right\}=\Omega
\]
\SourceLeaf{473}les forces \(\xi,\eta,\zeta\) sont exprimées par les différentielles partielles
\[
\xi=\mu\left(\frac{\mathrm{d}\Omega}{\mathrm{d}x}\right)\mathpunct{\raisebox{.45ex}{,}}\qquad
\eta=\mu\left(\frac{\mathrm{d}\Omega}{\mathrm{d}y}\right)\mathpunct{\raisebox{.45ex}{,}}\qquad
\zeta=\mu\left(\frac{\mathrm{d}\Omega}{\mathrm{d}z}\right)\mathpunct{\cdot}
\]
Ainsi \(\mathrm{d}\dfrac{2\mu}{a+\delta a}\) sera la différentielle de \(m\Omega\) en regardant comme constantes des quantités qui se rapportent à la planète perturbante. Maintenant il est clair que \(\Omega\) est une fonction \textit{périodique} de \(\boldsymbol{L}\) et de \(\boldsymbol{L}'\); en la développant donc en termes de la forme \(k\cos(iL+i'L')\) et \(k\sin(iL+i'L')\), sa différentielle prise en traitant \(L'\) comme constant ne contiendra aucun terme indépendant de \(L\). Et comme on a
\[
\begin{aligned}
n+\delta n&=\BareRoot{m}.(a+\delta a)^{-\frac{3}{2}}\\
\mathrm{d}(n+\delta n)&=-\tfrac{3}{2}\BareRoot{m}.(a+\delta a)^{-\frac{5}{2}}\mathrm{d}(a+\delta a)=\tfrac{3}{2}\BareRoot{\dfrac{m}{a+\delta a}}\cdot\mathrm{d}\frac{1}{a+\delta a}\mathpunct{\raisebox{.45ex}{,}}
\end{aligned}
\]
on aura, en n’ayant égard qu’aux termes du premier ordre,
\[
\mathrm{d}(n+\delta n)=3\frac{\mu}{\BareRoot{ma}}\cdot\mathrm{d}\Omega,
\]
la différentielle \(\mathrm{d}\Omega\) prise relativement à \(L\), d’où il s’ensuit que la variation \(\mathrm{d}(n+\delta n)\) a la forme indiquée ci-dessus.
\par
En supposant que \(\bar n\) est la constante introduite par l’intégration de \(\mathrm{d}(n+\delta n)\), il est clair que ce qui répond au mouvement moyen, c. à d. l’intégrale \(\int n\mathrm{d}t\) sera composée
\begin{enumerate}\setlength{\itemsep}{0pt}\setlength{\parsep}{0pt}
\item[1)] de termes périodiques
\item[2)] d’un terme proportionnel au\raisebox{.45ex}{-} tems, \(=\bar n t\)
\item[3)] d’une partie constante.
\end{enumerate}
\par
En réunissant donc ce mouvement moyen à l’époque variable, qui elle-même sera aussi composée
\begin{enumerate}\setlength{\itemsep}{0pt}\setlength{\parsep}{0pt}
\item[1)] de termes périodiques
\item[2)] d’une variation séculaire que nous désignerons par \(\bar\varepsilon t\)
\item[3)] d’une partie constante,
\end{enumerate}
\noindent la longitude moyenne aura la forme \(\bar L+(\bar n+\bar\varepsilon)t+P\), \(P\) désignant les parties périodiques.
\par
De cette manière il nous reste six termes constans, que nous nommerons les \textit{valeurs moyennes des élémens}. En les réunissant avec les parties \SourceLeaf{474}périodiques et les variations séculaires on aura les valeurs complètes des élémens pour tout instant.
\par
Il conviendra de choisir pour les parties constantes des élémens \(\AscendingNode,\ i,\ \varphi,\ \tilde\omega\), d’après lesquelles on calcule les perturbations, les valeurs moyennes mêmes, c’est à dire les constantes introduites par l’intégration, et pour [celles de] \(L\) et \(n\) les quantités \(\bar L\) et \(\bar n+\bar\varepsilon\). Mais on voit que pour effectuer cela dans la pratique un calcul double sera indispensable. En effet après avoir [fait] le premier calcul, selon des élémens approchés quelconques, on déterminera par des méthodes connues les valeurs qu’il faut donner aux six constantes pour obtenir le meilleur accord avec les observations, et on répétera le calcul une seconde fois, en partant de ces valeurs. À moins que les premières valeurs n’aient été trop loin de la vérité, on ne sera certainement pas dans le cas de faire ce calcul pour la troisième fois.
\par\endgroup
